\documentclass[11pt,a4paper]{article}
\usepackage[utf8]{inputenc}
\usepackage[T1]{fontenc}
\usepackage[spanish]{babel}
\usepackage{amsmath,amssymb}
\usepackage{tikz}
\usetikzlibrary{patterns,calc}
\usepackage[margin=2cm]{geometry}
\usepackage{enumitem}
\usepackage{booktabs}
\usepackage{multicol}
\usepackage{xcolor}

% Colores
\definecolor{azul}{RGB}{70,130,180}
\definecolor{rojo}{RGB}{178,34,34}
\definecolor{verde}{RGB}{34,139,34}
\definecolor{gris}{RGB}{128,128,128}
\definecolor{fondo}{RGB}{245,245,245}

\title{\textbf{Reglas y Propiedades de Simplificación de Conjuntos}\\[0.3em]
\large Material de Referencia -- MPI I}
\author{}
\date{}

\begin{document}
\maketitle
\thispagestyle{empty}

%% ============================================================
%% SECCIÓN 1: OPERACIONES BÁSICAS
%% ============================================================
\section{Operaciones Básicas}

\begin{minipage}[t]{0.48\textwidth}
\textbf{Unión ($A \cup B$)}

\begin{tikzpicture}[scale=0.7]
  \draw[thick] (0,0) circle (1.8cm);
  \draw[thick] (1.2,0) circle (1.8cm);
  \fill[pattern=north east lines, pattern color=azul!60] (-1.8,0) circle (1.8cm);
  \fill[pattern=north east lines, pattern color=azul!60] (1.2,0) circle (1.8cm);
  \draw[thick] (0,0) circle (1.8cm);
  \draw[thick] (1.2,0) circle (1.8cm);
  \node at (-0.9,0) {\Large $A$};
  \node at (2.1,0) {\Large $B$};
  \draw (-2.5,-2.3) rectangle (4,2.3);
  \node at (0.6,2.6) {\small $U$};
\end{tikzpicture}

$A \cup B = \{x \mid x \in A \lor x \in B\}$
\end{minipage}
\hfill
\begin{minipage}[t]{0.48\textwidth}
\textbf{Intersección ($A \cap B$)}

\begin{tikzpicture}[scale=0.7]
  \draw[thick] (0,0) circle (1.8cm);
  \draw[thick] (1.2,0) circle (1.8cm);
  \fill[pattern=north east lines, pattern color=rojo!60] (0.6,0) ellipse (0.6cm and 1.5cm);
  \draw[thick] (0,0) circle (1.8cm);
  \draw[thick] (1.2,0) circle (1.8cm);
  \node at (-0.9,0) {\Large $A$};
  \node at (2.1,0) {\Large $B$};
  \draw (-2.5,-2.3) rectangle (4,2.3);
  \node at (0.6,2.6) {\small $U$};
\end{tikzpicture}

$A \cap B = \{x \mid x \in A \land x \in B\}$
\end{minipage}

\vspace{1em}

\begin{minipage}[t]{0.48\textwidth}
\textbf{Diferencia ($A - B$)}

\begin{tikzpicture}[scale=0.7]
  \begin{scope}
    \clip (0,0) circle (1.8cm);
    \fill[pattern=north east lines, pattern color=verde!60] (-2.5,-2.3) rectangle (4,2.3);
  \end{scope}
  \begin{scope}
    \clip (1.2,0) circle (1.8cm);
    \fill[white] (-2.5,-2.3) rectangle (4,2.3);
  \end{scope}
  \draw[thick] (0,0) circle (1.8cm);
  \draw[thick] (1.2,0) circle (1.8cm);
  \node at (-0.9,0) {\Large $A$};
  \node at (2.1,0) {\Large $B$};
  \draw (-2.5,-2.3) rectangle (4,2.3);
  \node at (0.6,2.6) {\small $U$};
\end{tikzpicture}

$A - B = \{x \mid x \in A \land x \notin B\}$
\end{minipage}
\hfill
\begin{minipage}[t]{0.48\textwidth}
\textbf{Complemento ($A'$)}

\begin{tikzpicture}[scale=0.7]
  \fill[pattern=north east lines, pattern color=gris!60] (-2.5,-2.3) rectangle (4,2.3);
  \begin{scope}
    \clip (0,0) circle (1.8cm);
    \fill[white] (-2.5,-2.3) rectangle (4,2.3);
  \end{scope}
  \draw[thick] (0,0) circle (1.8cm);
  \draw[thick] (1.2,0) circle (1.8cm);
  \node at (-0.9,0) {\Large $A$};
  \node at (2.1,0) {\Large $B$};
  \draw (-2.5,-2.3) rectangle (4,2.3);
  \node at (0.6,2.6) {\small $U$};
\end{tikzpicture}

$A' = U - A = \{x \in U \mid x \notin A\}$
\end{minipage}

\vspace{1.5em}

%% ============================================================
%% SECCIÓN 2: PROPIEDADES FUNDAMENTALES
%% ============================================================
\section{Propiedades Fundamentales}

\begin{multicols}{2}
\begin{itemize}[leftmargin=*, itemsep=0.8em]
  \item \textbf{Conmutativa:}
    \begin{align*}
      A \cup B &= B \cup A \\
      A \cap B &= B \cap A
    \end{align*}

  \item \textbf{Asociativa:}
    \begin{align*}
      (A \cup B) \cup C &= A \cup (B \cup C) \\
      (A \cap B) \cap C &= A \cap (B \cap C)
    \end{align*}

  \item \textbf{Distributiva:}
    \begin{align*}
      A \cup (B \cap C) &= (A \cup B) \cap (A \cup C) \\
      A \cap (B \cup C) &= (A \cap B) \cup (A \cap C)
    \end{align*}

  \item \textbf{Idempotente:}
    \begin{align*}
      A \cup A &= A \\
      A \cap A &= A
    \end{align*}

  \item \textbf{Absorción:}
    \begin{align*}
      A \cup (A \cap B) &= A \\
      A \cap (A \cup B) &= A
    \end{align*}
\end{itemize}
\end{multicols}

\vspace{1em}

%% ============================================================
%% SECCIÓN 3: ELEMENTOS NEUTROS Y ABSORBENTES
%% ============================================================
\section{Elementos Neutros y Absorbentes}

\begin{center}
\begin{tabular}{lccc}
\toprule
\textbf{Tipo} & \textbf{Operación} & \textbf{Resultado} & \textbf{Interpretación} \\
\midrule
Neutro & $A \cup \varnothing$ & $= A$ & El vacío no aporta nada \\
Neutro & $A \cap U$ & $= A$ & El universal no cambia nada \\
Absorbente & $A \cup U$ & $= U$ & Todo se ``absorbe'' en $U$ \\
Absorbente & $A \cap \varnothing$ & $= \varnothing$ & Todo se ``absorbe'' en $\varnothing$ \\
\bottomrule
\end{tabular}
\end{center}

\vspace{1em}

%% ============================================================
%% SECCIÓN 4: COMPLEMENTO
%% ============================================================
\section{Complemento}

\begin{multicols}{2}
\begin{itemize}[leftmargin=*, itemsep=0.8em]
  \item \textbf{Doble complemento:} $(A')' = A$
  \item $\varnothing' = U$
  \item $U' = \varnothing$
  \item \textbf{Complemento propio:}
    \begin{align*}
      A \cup A' &= U \\
      A \cap A' &= \varnothing
    \end{align*}
\end{itemize}
\end{multicols}

\vspace{1em}

%% ============================================================
%% SECCIÓN 5: LEYES DE DE MORGAN
%% ============================================================
\section{Leyes de De Morgan -- \textcolor{rojo}{CRÍTICO}}

\begin{center}
\fbox{\parbox{0.85\textwidth}{
\begin{center}
\Large
$\overline{A \cup B} = \overline{A} \cap \overline{B}$

\vspace{0.5em}
\small El complemento de una unión $=$ la intersección de los complementos

\vspace{1em}
$\overline{A \cap B} = \overline{A} \cup \overline{B}$

\vspace{0.5em}
\small El complemento de una intersección $=$ la unión de los complementos
\end{center}
}}
\end{center}

\vspace{0.5em}

\textbf{Truco mnemotécnico:} Al pasar el complemento, se invierte $\cup \leftrightarrow \cap$

\vspace{0.5em}

\begin{minipage}[t]{0.48\textwidth}
\textbf{$\overline{A \cup B}$}

\begin{tikzpicture}[scale=0.7]
  \fill[pattern=north east lines, pattern color=gris!60] (-2.5,-2.3) rectangle (4,2.3);
  \begin{scope}
    \clip (0,0) circle (1.8cm);
    \fill[white] (-2.5,-2.3) rectangle (4,2.3);
  \end{scope}
  \begin{scope}
    \clip (1.2,0) circle (1.8cm);
    \fill[white] (-2.5,-2.3) rectangle (4,2.3);
  \end{scope}
  \draw[thick] (0,0) circle (1.8cm);
  \draw[thick] (1.2,0) circle (1.8cm);
  \node at (-0.9,0) {\Large $A$};
  \node at (2.1,0) {\Large $B$};
  \draw (-2.5,-2.3) rectangle (4,2.3);
  \node at (0.6,2.6) {\small $U$};
\end{tikzpicture}

$\overline{A \cup B} = \overline{A} \cap \overline{B}$
\end{minipage}
\hfill
\begin{minipage}[t]{0.48\textwidth}
\textbf{$\overline{A \cap B}$}

\begin{tikzpicture}[scale=0.7]
  \fill[pattern=north east lines, pattern color=gris!60] (-2.5,-2.3) rectangle (4,2.3);
  \begin{scope}
    \clip (0.6,0) ellipse (0.6cm and 1.5cm);
    \fill[white] (-2.5,-2.3) rectangle (4,2.3);
  \end{scope}
  \draw[thick] (0,0) circle (1.8cm);
  \draw[thick] (1.2,0) circle (1.8cm);
  \node at (-0.9,0) {\Large $A$};
  \node at (2.1,0) {\Large $B$};
  \draw (-2.5,-2.3) rectangle (4,2.3);
  \node at (0.6,2.6) {\small $U$};
\end{tikzpicture}

$\overline{A \cap B} = \overline{A} \cup \overline{B}$
\end{minipage}

\vspace{1.5em}

%% ============================================================
%% SECCIÓN 6: DIFERENCIA Y COMPLEMENTO
%% ============================================================
\section{Diferencia y Complemento}

\begin{multicols}{2}
\begin{itemize}[leftmargin=*, itemsep=0.8em]
  \item $A - B = A \cap B'$
  \item $A - B = \overline{A' \cup B}$
  \item $A \triangle B = (A - B) \cup (B - A)$ \hfill\textit{(diferencia simétrica)}
  \item $A \triangle B = (A \cup B) - (A \cap B)$
\end{itemize}
\end{multicols}

\vspace{0.5em}

\begin{minipage}[t]{0.48\textwidth}
\textbf{Diferencia Simétrica ($A \triangle B$)}

\begin{tikzpicture}[scale=0.7]
  \begin{scope}
    \clip (0,0) circle (1.8cm);
    \fill[pattern=north east lines, pattern color=verde!60] (-2.5,-2.3) rectangle (4,2.3);
  \end{scope}
  \begin{scope}
    \clip (1.2,0) circle (1.8cm);
    \fill[pattern=north east lines, pattern color=azul!60] (-2.5,-2.3) rectangle (4,2.3);
  \end{scope}
  \begin{scope}
    \clip (0.6,0) ellipse (0.6cm and 1.5cm);
    \fill[white] (-2.5,-2.3) rectangle (4,2.3);
  \end{scope}
  \draw[thick] (0,0) circle (1.8cm);
  \draw[thick] (1.2,0) circle (1.8cm);
  \node at (-0.9,0) {\Large $A$};
  \node at (2.1,0) {\Large $B$};
  \draw (-2.5,-2.3) rectangle (4,2.3);
  \node at (0.6,2.6) {\small $U$};
\end{tikzpicture}

$A \triangle B = (A \cup B) - (A \cap B)$
\end{minipage}
\hfill
\begin{minipage}[t]{0.48\textwidth}
\vspace{2.5cm}
\end{minipage}

\vspace{1em}

%% ============================================================
%% SECCIÓN 7: TABLA DE SIMPLIFICACIÓN
%% ============================================================
\section{Tabla de Simplificación Rápida}

\begin{center}
\begin{tabular}{lcc}
\toprule
\textbf{Expresión} & \textbf{Se simplifica a} & \textbf{Ley} \\
\midrule
$A \cup A$ & $A$ & Idempotente \\
$A \cap A$ & $A$ & Idempotente \\
$A \cup \varnothing$ & $A$ & Neutro \\
$A \cap \varnothing$ & $\varnothing$ & Absorbente \\
$A \cup U$ & $U$ & Absorbente \\
$A \cap U$ & $A$ & Neutro \\
$A \cup A'$ & $U$ & Complemento \\
$A \cap A'$ & $\varnothing$ & Complemento \\
$A \cup (A \cap B)$ & $A$ & Absorción \\
$A \cap (A \cup B)$ & $A$ & Absorción \\
$(A')'$ & $A$ & Dobble compl. \\
\bottomrule
\end{tabular}
\end{center}

\vspace{1em}

%% ============================================================
%% SECCIÓN 8: EQUIVALENCIAS ÚTILES
%% ============================================================
\section{Equivalencias Útiles para Simplificar}

\begin{center}
\begin{tabular}{lcl}
\toprule
\textbf{Para simplificar} & & \textbf{Reemplazar por} \\
\midrule
$A - B$ & $\longleftrightarrow$ & $A \cap B'$ \\
$\overline{A \cup B}$ & $\longleftrightarrow$ & $\overline{A} \cap \overline{B}$ \\
$\overline{A \cap B}$ & $\longleftrightarrow$ & $\overline{A} \cup \overline{B}$ \\
$A \subseteq B$ & $\longleftrightarrow$ & $A \cap B = A$ \\
$A \subseteq B$ & $\longleftrightarrow$ & $A \cup B = B$ \\
\bottomrule
\end{tabular}
\end{center}

\vspace{1em}

%% ============================================================
%% SECCIÓN 9: RELACIONES ENTRE ⊆ Y OPERACIONES
%% ============================================================
\section{Relaciones entre $\subseteq$ y Operaciones}

\begin{center}
\fbox{\parbox{0.8\textwidth}{
\begin{center}
$A \subseteq B \iff A \cap B = A \iff A \cup B = B \iff A - B = \varnothing \iff B' \subseteq A'$
\end{center}
}}
\end{center}

\vspace{1em}

%% ============================================================
%% SECCIÓN 10: TRUCOS DE SIMPLIFICACIÓN
%% ============================================================
\section{Trucos de Simplificación Paso a Paso}

\begin{enumerate}[leftmargin=*, itemsep=0.5em]
  \item \textbf{Reemplazar diferencias:} $A - B \to A \cap B'$
  \item \textbf{Aplicar De Morgan:} sacar complementos de expresiones compuestas
  \item \textbf{Buscar absorción:} $A \cap (A \cup B) \to A$
  \item \textbf{Buscar complementos:} $A \cap A' \to \varnothing$
  \item \textbf{Buscar idempotencia:} $A \cup A \to A$
  \item \textbf{Repetir} hasta que no haya más simplificaciones
\end{enumerate}

\vspace{1em}

\textbf{Ejemplo resuelto:}

\begin{center}
\fbox{\parbox{0.9\textwidth}{
\textbf{Simplificar:} $(A \cap B) \cup (A \cap B')$

\vspace{0.5em}
\textbf{Paso 1:} Factorizar $A$: $A \cap (B \cup B')$

\vspace{0.5em}
\textbf{Paso 2:} $B \cup B' = U$ (complemento propio)

\vspace{0.5em}
\textbf{Paso 3:} $A \cap U = A$ (neutro)

\vspace{0.5em}
\textbf{Resultado:} $\boxed{A}$
}}
\end{center}

\vfill
\begin{center}
\rule{0.8\textwidth}{0.4pt}

\vspace{0.5em}
Generado en \LaTeX{} por: \textbf{Jonathan Javier Pietracci}
\end{center}

\end{document}
